\textbf{Background.} Published accuracies for radiograph body-part classification are hard to compare: splits, thresholds and uncertainty are rarely reported, and most large studies use private data.
\textbf{Methods.} We built a leakage-controlled benchmark on the public UNIFESP dataset (22 classes, 1,738 radiographs, 768 patients) with patient-level folds, a 323-image held-out set, all choices made on out-of-fold predictions and patient-bootstrap confidence intervals. We added ablations, a frozen foundation-model baseline, transfer to MURA and FracAtlas, a MURA benchmark, calibration and selective prediction, and a label audit.
\textbf{Results.} The ensemble reached a held-out sample-mean F1 of 0.919 (95 \% CI 0.867 to 0.960), not significantly above its best single model (+0.0066, p = 0.063); no ablation differed from the reference (all p > 0.08). A gated CNN-Transformer fusion ignored its CNN branch. A random image-level split overstated validation F1 by 0.0144 (patient-level folds: 0.0004), and tuning thresholds on the scored set added 0.0083. On the official MURA validation set accuracy was 0.9747 (0.9663 to 0.9818), against 0.9643 reported by Shah et al. for the same images. The ensemble was well calibrated (ECE 0.054), and accuracy on retained images rose from 0.926 to 0.945 at 90 \% coverage. The audit flagged 39 of 1,415 training images (2.8 \%) as possible label errors, unreviewed by clinicians. A late Kaggle submission without pseudo-labels scored 0.9256 (private), between second and third place.
\textbf{Conclusion.} How a classifier is evaluated matters more than its architecture: gains from fusion, ensembling and pseudo-labelling fell within the noise of a 323-image held-out set.


